% ============================================================================
% chap5_dual_visualization.tex
%   第5章 §5.2.2 · 双通道可视化机制
%   语义级 WebGUI 通道 + 原生 noVNC 通道 的并行数据链路
% ============================================================================

\documentclass[tikz,border=6pt]{standalone}
\usepackage[UTF8]{ctex}
\usepackage{amsmath}
\usepackage{tikz}
\usetikzlibrary{positioning, arrows.meta, shapes.geometric, calc, fit, backgrounds}

% 颜色：边框和箭头颜色保持不变，文字统一改为黑色
\definecolor{cOurs}{HTML}{0B4F8A}      % 深蓝：原生通道
\definecolor{cOursAlt}{HTML}{C05600}   % 深橙：源头/前端
\definecolor{cOursEval}{HTML}{000000}  % 黑色：语义通道
\definecolor{cUpstream}{HTML}{4D4D4D}  % 深灰：通用基础模块
\definecolor{cBgLayer}{HTML}{F0F0F0}
\definecolor{cBgHi}{HTML}{FFF8E7}

\begin{document}
\begin{tikzpicture}[
  font=\large\sffamily,
  every node/.style={align=center, text=black},
  % 源/宿 (左右两端)
  endpoint/.style={
    draw=cOursAlt, line width=1.1pt, rounded corners=3pt,
    fill=cBgHi, text=black, inner sep=4pt,
    minimum width=3.2cm, minimum height=2.1cm,
    font=\normalsize\bfseries\sffamily,
  },
  % 语义通道节点
  semnode/.style={
    draw=cOursEval, line width=1pt, rounded corners=3pt,
    fill=white, text=black, inner sep=4pt,
    minimum width=2.6cm, minimum height=1.3cm,
    font=\normalsize\sffamily,
  },
  % 原生通道节点
  vncnode/.style={
    draw=cOurs, line width=1pt, rounded corners=3pt,
    fill=white, text=black, inner sep=4pt,
    minimum width=2.6cm, minimum height=1.3cm,
    font=\normalsize\sffamily,
  },
  % 通道标题横幅
  banner/.style={
    font=\normalsize\bfseries\sffamily,
    inner sep=3pt, rounded corners=2pt,
    text=black,
  },
  % 箭头
  semflow/.style={
    -{Stealth[length=2.4mm, width=2mm]},
    line width=0.9pt, draw=cOursEval,
  },
  vncflow/.style={
    -{Stealth[length=2.4mm, width=2mm]},
    line width=0.9pt, draw=cOurs,
  },
  ctrlflow/.style={
    -{Stealth[length=2.2mm, width=1.8mm]},
    line width=0.9pt, draw=black!75, densely dashed,
  },
  lbl/.style={font=\normalsize, inner sep=1pt, fill=white, text=black},
]

% ============================================================================
% 左端: 仿真源
% ============================================================================
\node[endpoint] (sim) at (-8.0cm, 0)
  {仿真器 (容器内)\\[2pt]
   \small Webots / Gazebo / RViz\\
   \small ROS\,2 状态话题};

% ============================================================================
% 右端: 浏览器前端
% ============================================================================
\node[endpoint] (browser) at (8.0cm, 0)
  {浏览器前端\\[2pt]
   \small React SPA\\
   \small WebGUI + iframe};

% ============================================================================
% 上行: 语义级 WebGUI 通道
% ============================================================================
\node[banner, text=black, fill=cOursEval!10, draw=cOursEval!60]
  (semtitle) at (0, 3.00cm)
  {语义级 WebGUI 通道 —— 结构化状态 \& 抽象图形};

\node[semnode] (extract) at (-4.3cm, 1.75cm)
  {状态提取器\\
   \small 位姿 / 扫描\\[-1pt]
   \small 安全集};

\node[semnode] (ws) at (0, 1.75cm)
  {WebSocket 推送\\
   \small JSON 数据\\[-1pt]
   \small 组件刷新};

\node[semnode] (comp) at (4.3cm, 1.75cm)
  {语义组件层\\
   \small 轨迹 / 力向量\\[-1pt]
   \small 约束边界};

\draw[semflow] (sim.north east) to[out=20,in=180] (extract.west);
\draw[semflow, shorten <=5pt, shorten >=5pt] (extract.east) -- (ws.west);
\draw[semflow, shorten <=5pt, shorten >=5pt] (ws.east) -- (comp.west);
\draw[semflow] (comp.east) to[out=0,in=160] (browser.north west);

% ============================================================================
% 下行: 原生 noVNC 通道
% ============================================================================
\node[banner, text=black, fill=cOurs!10, draw=cOurs!60]
  (vnctitle) at (0, -3.00cm)
  {原生 noVNC 通道 —— 仿真器 GUI 远程转发};

\node[vncnode] (xvfb) at (-4.3cm, -1.75cm)
  {Xvfb 虚拟显示\\
   \small 无头 X server\\[-1pt]
   \small GUI 承载};

\node[vncnode] (xvnc) at (0, -1.75cm)
  {Xvnc 服务\\
   \small VNC 帧编码\\[-1pt]
   \small 端口转发};

\node[vncnode] (novnc) at (4.3cm, -1.75cm)
  {noVNC 客户端\\
   \small WebSocket 解码\\[-1pt]
   \small iframe 嵌入};

\draw[vncflow] (sim.south east) to[out=-20,in=180] (xvfb.west);
\draw[vncflow, shorten <=5pt, shorten >=5pt] (xvfb.east) -- (xvnc.west);
\draw[vncflow, shorten <=5pt, shorten >=5pt] (xvnc.east) -- (novnc.west);
\draw[vncflow] (novnc.east) to[out=0,in=-160] (browser.south west);

% ============================================================================
% 控制反向链路
% ============================================================================
\draw[ctrlflow] (browser.west) to[out=175,in=5]
  node[lbl, pos=0.5]{用户控制 (start / pause / reset)}
  (sim.east);

% ============================================================================
% 图例
% ============================================================================
\node[draw=black!30, rounded corners=2pt, inner sep=5pt, font=\normalsize,
      text=black, anchor=north west]
  at (3.8cm, -3.5cm)
  {\begin{tabular}{@{}l@{\hspace{4pt}}l@{}}
     \tikz[baseline=-0.6ex]\draw[semflow] (0,0) -- (0.7,0); & 语义数据流 (JSON) \\[2pt]
     \tikz[baseline=-0.6ex]\draw[vncflow] (0,0) -- (0.7,0); & 原生帧流 (VNC) \\[2pt]
     \tikz[baseline=-0.6ex]\draw[ctrlflow] (0,0) -- (0.7,0); & 用户控制指令 \\
   \end{tabular}};

\end{tikzpicture}
\end{document}